\documentclass[UTF8,11pt]{ctexart}
\usepackage[a4paper,margin=24mm]{geometry}
\usepackage{amsmath,amssymb,amsthm,mathtools}
\usepackage[all]{xy}
\usepackage{xurl,hyperref,enumitem}
\hypersetup{hidelinks,breaklinks=true}
\setlength{\emergencystretch}{3em}
\newtheorem*{theorem}{Theorem}
\newtheorem*{lemma}{Lemma}
\newtheorem*{proposition}{Proposition}
\newtheorem*{corollary}{Corollary}
\theoremstyle{definition}
\newtheorem*{definition}{Definition}
\newtheorem*{remark}{Remark}
\DeclareMathOperator{\Hom}{Hom}
\DeclareMathOperator{\Ext}{Ext}
\DeclareMathOperator{\Tor}{Tor}
\DeclareMathOperator{\Spec}{Spec}
\DeclareMathOperator{\supp}{supp}
\DeclareMathOperator{\im}{im}
\DeclareMathOperator{\id}{id}
\DeclareMathOperator{\Tr}{Tr}
\DeclareMathOperator{\rank}{rank}
\newcommand{\R}{\mathbb R}
\newcommand{\C}{\mathbb C}
\newcommand{\Z}{\mathbb Z}
\newcommand{\N}{\mathbb N}
\newcommand{\Q}{\mathbb Q}
\begin{document}
\section*{011\quad 有限步增边下随机游走类型不变}
\subsection*{来源与收录范围}
Peter G. Doyle and J. Laurie Snell, \emph{Random Walks and Electric Networks}, version dated 5 July 2006. 主命题为 \S2.4.4 的未编号 Theorem，印刷 pp.~103--104；并收入 \S2.4.3 的电导比较定理及其证明（pp.~102--103），作为本题支撑，不另计题数。
\par\url{https://math.dartmouth.edu/~doyle/docs/walks/walks.pdf}
\par 获取日期：2026-09-28。正文据已取得的 PDF 文本转录；第103页已查看页面图像，第104页截图接口失败。未取得可保存的原始 PDF；包内转录摘录不冒称原件。
\paragraph{本题使用的上下文（据原书 \S2.4.1--2.4.2 整理，非作者逐字原句）}
$G$ 为无限连通有界度图，顶点度数至多为整数 $E$，$k$ 为正整数。电网络 $(G,C)$ 的转移概率为 $P_{xy}=C_{xy}/C_x$，$C_x=\sum_y C_{xy}$。从参考顶点 $0$ 出发不再返回的概率为零称 recurrent，否则称 transient；这两个可能性称为 type。$G_k$ 为在图距离至多 $k$ 的顶点之间增边得到的 $k$-fuzz。主目标是证明增边前后的简单随机游走具有相同 type，而不是仅证明某个固定图的常返性。
\subsection*{作者命题与人类证明}
\paragraph{Comparing two networks (Section 2.4.3).}
Given two sets of conductances $C$ and $\bar C$ on $G$, we say that $(G,\bar C)<(G,C)$ if $\bar C_{xy}<C_{xy}$ for all $xy$, or equivalently, if $\bar R_{xy}>R_{xy}$ for all $xy$. Assume that $(G,\bar C)<(G,C)$. Then by the Monotonicity Law, $\bar R_{\mathrm{eff}}\ge R_{\mathrm{eff}}$. Thus if random walk on $(G,\bar C)$ is transient, i.e., if $\bar R_{\mathrm{eff}}<\infty$, then random walk on $(G,C)$ is also transient. If random walk on $(G,C)$ is recurrent, i.e., if $R_{\mathrm{eff}}=\infty$, then random walk on $(G,\bar C)$ is also recurrent.
\begin{theorem}[Section 2.4.3]
If $(G,C)$ and $(G,\bar C)$ are networks, and if there exist constants $u,v$ with $0<u\le v<\infty$ such that
\[
uC_{xy}\le\bar C_{xy}\le vC_{xy}
\]
for all $x$ and $y$, then random walk on $(G,\bar C)$ is of the same type as random walk on $(G,C)$.
\end{theorem}
\begin{proof}
Let $U_{xy}=uC_{xy}$ and $V_{xy}=vC_{xy}$. Then $(G,U)\le(G,\bar C)\le(G,V)$. But the random walks for $(G,U)$ and $(G,V)$ are the same as random walk on $(G,C)$. Thus random walk for $(G,\bar C)$ is of the same type as random walk on $(G,C)$.
\end{proof}
\begin{corollary}[Section 2.4.3]
Let $(G,C)$ be a network. If for every edge $xy$ of $G$ we have $0<u<C_{xy}<v<\infty$ for some constants $u$ and $v$, then the random walk on $(G,C)$ has the same type as simple random walk on $G$.
\end{corollary}
\paragraph{The $k$-fuzz of a graph (Section 2.4.4).}
\begin{theorem}
Simple random walk on $G$ and on the $k$-fuzz $G_k$ of $G$ have the same type.
\end{theorem}
\begin{proof}
Let $\mathbf P$ be simple random walk on $G$. Define $\bar{\mathbf P}=(\mathbf P+\mathbf P^2+\ldots+\mathbf P^k)/k$. Then $\bar{\mathbf P}$ may be considered to be $\mathbf P$, watched at one of the first $k$ steps chosen at random, then at a time chosen at random from the next $k$ steps after this time, etc. Thinking of $\bar{\mathbf P}$ in this way, we see that $\mathbf P$ is in state $0$ at least once for every time $\bar{\mathbf P}$ in state $0$. Hence, if $\bar{\mathbf P}$ is recurrent so is $\mathbf P$. Assume now that $\bar{\mathbf P}$ is transient. Choose a finite set $S$ so that $0$ cannot be reached in $k$ steps from a point outside of $S$. Then, since the walk $\bar{\mathbf P}$ will be outside $S$ from some time on, the walk $\mathbf P$ cannot be at $0$ after this time, and $\mathbf P$ is also transient. Therefore, $\mathbf P$ and $\bar{\mathbf P}$ are of the same type.

Finally, we show that $\bar{\mathbf P}$ has the same type as simple random walk on $G_k$. Here it is important to remember our restriction that $G$ is of bounded degree, so that for some $E$ no vertex has degree $>E$. We know that $\mathbf P$ is reversible with $\mathbf w\mathbf P=\mathbf w$, where $w_x$ is the number of edges coming out of $x$. From its construction, $\bar{\mathbf P}$ is also reversible and $\mathbf w\bar{\mathbf P}=\mathbf w$. $\bar{\mathbf P}$ is the random walk on a network $(G_k,\bar C)$ with $\bar C_{xy}=w_x\bar P_{xy}$. If $\bar P_{xy}>0$, there is a path $x,x_1,x_2,\ldots,x_{m-1},y$ in $G$ from $x$ to $y$ of length $m\le k$. Then
\[
\bar P_{xy}\ge\frac1k\left(\frac1E\right)^m
\ge\frac1k\left(\frac1E\right)^k.
\]
Thus
\[
0<\frac1k\left(\frac1E\right)^k\le\bar P_{xy}\le1
\]
and
\[
0<\frac1k\left(\frac1E\right)^k\le\bar C_{xy}\le E.
\]
Therefore, by the theorem on the irrelevance of bounded twiddling proven in Section 2.4.3, $\bar{\mathbf P}$ and simple random walk on $G_k$ are of the same type. So $G$ and $G_k$ are of the same type.
\end{proof}
\subsection*{整理说明（非作者正文）}
本条只计 $k$-fuzz 的类型不变性一个问题，不把电导比较定理或其推论另计题数。标准先修为不可约可数状态 Markov 链的常返/暂留分类、可逆链与电导的对应，以及无穷网络的有效电阻判据和 Rayleigh 单调性（原书 \S2.4.2--2.4.3 所指内容）。本题仍需完成随机观察时间对应、有限步转移的电导表示，以及由有界度导出的统一双边比较，不是直接应用某个已给的类型不变性定理。

原主证明的两个段落及全部界均已收录，保留原语言和顺序。省略的 Figure~61 是二维 $2$-fuzz 的例图，不承担一般定理证明；后续应用与习题不作为本题续证。
\subsection*{可修改的简短标注（非作者正文）}
$c$：无限图随机游走的常返与暂留类型

$a$：随机观察时间；可逆Markov链；转移概率与电导的对应；Rayleigh单调性；有界度与电导一致比较

\texttt{cross\_theory\_status = yes}。把有限步增边转成原链若干步转移的平均，再表示为电网络；有界度给出电导比较，返回随机游走类型不变。位置：\S2.4.4 全证，pp.~103--104。

全部标注均为非穷尽、可修改初稿；未标注不表示未使用。
\subsection*{许可与修改记录}
原数学正文作者为 Peter G. Doyle 与 J. Laurie Snell；原书 2006 年版本按 GNU Free Documentation License 发布，无不变章节、封面文字或封底文字。许可全文随包保留在 \texttt{sources/stacks/GFDL-1.2.txt}。本题标题不同于原书，整理部分随原文许可；不暗示原作者认可本次标注。2026-09-28：助手转录、加入来源定位和说明；未新增数学证明。
\end{document}
